\section{Trajectory overlays and row-column-row mixing}

\label{special:overlay}

We add independent switches along selected trajectories of a base shuffle. A character estimate mixes the resulting permutation on one group of indices. A sequence of calls on two transverse partitions then reduces full mixing to a local limit theorem for finite contingency tables.

Unlike the first route, which conditions on specified outside labels, this route conditions on a complete base shuffle. Independent extra switches then permute chosen groups of base indices. Predictability of their placement ensures that the resulting physical shuffle still has the original law. The useful output is uniformity of an entire group permutation conditional on the base, on average over that base.

Use $V=\mathbb F_2^d$, with the convention that the later permutation acts on the left. For a finite set $A$, write $\mathfrak S(A)$ for its symmetric group. As in Section~\ref{sec:conditional}, the layer directions are cyclic and their current phase may be arbitrary.

The \textbf{base shuffle} has exactly these layer dynamics. We first prove a partial mixing estimate to prepare the environments for the extra switches.

\begin{theorem}[Trajectory-overlay construction]\label{thm:overlay}
For $n=2^d$, the full-deck law from any deterministic start has total-variation distance tending to zero from uniform after $33030144d$ physical shuffles.
\end{theorem}

\subsection{Mixing a fraction of the base indices}

Fix \(K=8\), and let \(q=n/K\); throughout this section only sufficiently large \(d\) need be considered.

\begin{lemma}[Preparing a group of base indices]\label{lem:overlay-preparation}
After \(20d\) layers, the ordered positions of any \(q\) specified cards in the base shuffle have total-variation distance \(o(1)\) from uniform ordered distinct positions. The estimate is uniform in their starting positions and in the starting cyclic phase.
\end{lemma}
\begin{proof}
The same deferred-column realization appears in the companion on random coordinate frames~\cite[Lemma ``Independent shear realization'']{random-frames}; we give the construction here for the required one-eighth-deck estimate.

We introduce a randomized linear change of frame, just for proving this estimate. Restart notation at layer 1 for the layers in the estimate, with cyclically ordered \(i_1,i_2,\ldots\). Represent the layer permutation as \(L_t\widetilde T_t\), where \(\widetilde T_t\) is an independent ordinary fair layer along \(e_{i_t}\), and
\[
 L_t x=x+\ell_t(x)e_{i_t}.
\]
Here \(\ell_t\) are independent random linear forms subject to \(\ell_t(e_{i_t})=0\), uniform under that condition and independent of the \(\widetilde T_t\)'s. Each \(L_t\) consists of switches on the same matching, so this factorization has the required shuffle law. With \(A_t=L_t\cdots L_1\), \(A_0=I\), map locations after layer \(t\) through \(A_t^{-1}\). In this frame the layer dynamics conditional on the \(L_s\)'s are independent arrays of fair switches in directions
\[
 v_t=A_{t-1}^{-1}e_{i_t},
\]
by conjugating \(\widetilde T_t\). Their conditional laws depend only on the directions.

For \(t>d\), conditional on previous directions, \(v_t\) is uniform outside the hyperplane spanned by \(v_{t-d+1},\ldots,v_{t-1}\). In fact the columns of the inverse transform evolve by
\[
 A_s^{-1} e_j=A_{s-1}^{-1}e_j+\ell_s(e_j)v_s,
\]
using \(L_s^{-1}=L_s\). Thus
\[
 v_t=v_{t-d}+\sum_{s=t-d+1}^{t-1}\ell_s(e_{i_t})v_s .
\]
These coefficients are independent fair bits independent of all \(v_p\) for \(p<t\): in the indicated range they update only the column \(i_t\), not used as a direction again until time \(t\) (so they do not affect updates via \(v_s\) before then either). All coordinates of each form except its constrained one are independent on the coordinate basis. The first \(d\) directions are independent in the linear-algebra sense, by triangularity of the updates, and inductively every window of \(d\) directions is linearly independent by the last display. This proves the assertion about the hyperplane.

Order the \(q\) cards. In the changed frame, let \(p_t\) be the conditional position distribution of card \(j\) given the directions through layer \(t\) and the paths through layer \(t\) of cards \(1,\ldots,j-1\). Write \(F_t\) for the free positions, i.e. those not occupied then by the given \(j-1\) cards, with size \(r=n-j+1\). We track \(u_t(x)=p_t(x)-1/r\) on \(F_t\). The recursive conditional-law calculation can be done as follows. Reveal the new direction, then the updates of the given path cards. Neither revelation biases the previous law of card \(j\) given the history: the new direction's law given past directions is independent of the past switches, and the updates of the path cards in this direction are from fair switches on their edges with law independent of card \(j\)'s location. Conditional on the updates, the switches on the other edges are still fair and independent. On any pair of two free sites the probabilities of card \(j\) are averaged. For a free site paired with a path card, its probability is deterministically moved to the new free site of that pair. These operations carry the uniform law on the free sites to the new such uniform law. Consequently \(\sum u_t^2\) does not increase in an update, and its decrease includes \((u_t(x)-u_t(y))^2/2\) for each matched pair of free sites before the update.

After the first \(d\) layers, split the free sites just before an update into \(F^0,F^1\) using the two cosets of the hyperplane for its direction. Given the history so far, every opposite-side pair is matched with probability \(2/n\). Both sides have size at least \(n/2-q\). Suppressing the time subscript, we have
\[
 \sum_{x\in F^0,y\in F^1}(u(x)-u(y))^2\ \ge\ (n/2-q)\sum_x u(x)^2
\]
because the two sums of \(u\) on the respective sides are negatives of one another. The conditional expected ratio bound for the new squared norm relative to the old is therefore \(1/2+q/n\), meaning the new expectation is at most this multiple of the old squared norm. Since the initial squared norm is at most 1, at \(T=20d\) the expected total-variation distance of the conditional law from uniform on the free sites is at most
\[
 \tfrac12\sqrt{n}\,(1/2+q/n)^{(T-d)/2}
\]
by Cauchy-Schwarz.

At this last time the conditional law given the full linear transforms and the path data in their frame only needs the directions and the path data, by the conditional switch laws. Undoing the final linear transform preserves the comparison with the uniform free-site law. Giving only the linear transforms and the previous cards' \emph{final positions} instead of their paths cannot increase the expected bound (by mixing the conditional laws; the uniform comparison depends only on final positions given the transforms). Summing the bound over \(j\le q\) bounds the expectation over the transforms of the distance of the ordered \(q\)-tuple law conditional on them from uniform. Here we used the sequential total-variation bound by the sum of expected distances of the successive conditional laws from the uniform remaining-site laws; it follows by replacing conditionals from last to first and using the triangle inequality. Our sum tends to zero since \(q=n/8\) and \(n=2^d\). Unconditioning proves the partial mixing estimate.
\end{proof}

\subsection{An overlay on the base trajectories}

The marginal estimate has prepared the starting positions of the indices that will receive extra switches. We now define those switches so that their conditional law can be studied separately while the resulting physical process remains an ordinary Thorp shuffle.

Call the initial position names of the base shuffle its \textbf{indices}, in \(V\). Just before a layer, the matching of locations induces a matching of the indices via their base positions at that time. We can put extra switches on selected edges of this index matching chosen from the pre-layer base data. In our overlay construction we supply these switches from fresh arrays of independent fair coins independent of the base arrays. Form the actual layer coins on the location pairs by xoring the corresponding extra bits (0 where absent) with the base bits. The actual shuffle has the ordinary law: conditional on the past and the extra bits for a layer, the base bits in that layer still have their independent uniform law, since the selection and extra bits do not use the current base coins. So the resulting coin array has the independent fair law at every layer given the past. If the base permutation (indices to current positions) is \(B\), the actual permutation is \(B D\), with \(D\) the overlay permutation on indices, initialized to the identity. At each layer \(D\) is left-multiplied by the extra switches in the pre-layer index matching and \(B\) is then updated by the base switches on locations. This gives exactly the xor dynamics just specified.

For a fixed group \(A\subseteq V\) of \(q\) indices, we call the following procedure a \emph{call} on \(A\).

\begin{proposition}[A conditional group permutation]\label{prop:overlay-call}
Let $A\subset V$ have $q=n/8$ indices. Use $k=65536$ phases, each consisting of $20d$ preparation layers without extra switches and one active sweep with extra switches on the pairs whose two base indices belong to $A$. Let $\kappa_{\mathcal B}$ be the law of the resulting multiplier on $\mathfrak S(A)$ conditional on the full base shuffle $\mathcal B$. Then
\[
 \sup_{\text{pre-call base histories}}
 \E\big[\|\kappa_{\mathcal B}-U_A\|_{\TV}
          \mid\text{pre-call history}\big]\longrightarrow0.
\]
The assertion is uniform in the starting cyclic phase and in the initial base placement. The multiplier law is independent of the incoming overlay permutation.
\end{proposition}

Given the base trajectories, the left multiplier is a product of independent random permutations on \(A\), with laws determined, phase by phase, by the segments of ordered base trajectories of \(A\) in the active layers (including their positions just before the active layers). These laws describe the multipliers irrespective of the incoming overlay. For analyzing the expected conditional distance we can use an ideal experiment with \emph{independent} such segments, each started from uniform ordered distinct locations of the indices in \(A\) and run with ordinary base layers. Indeed the joint law of the actual segments, even given data before the call, is within \(o(1)\) in total variation of this ideal segment law for constant \(k\). This follows by Lemma~\ref{lem:overlay-preparation} in each \(20d\)-layer preparation interval, conditional on earlier base history, followed by the same ordered-location transition dynamics on the active layers. These dynamics do not depend on the placements of other indices. Equivalently one can couple each new segment start to a fresh uniform one and use matched segment dynamics when the coupling succeeds, with failure probability bounded by the sum of the partial mixing errors. Since the conditional distance is a bounded function of the segments, it suffices to bound its expectation in the ideal experiment.

It remains to prove the conditional group estimate for these independent ideal segments, which we do next.

\subsection{Estimate for an ideal active segment}\label{sec:overlay-ideal}

Let \(E\) denote an ideal segment environment, i.e. the ordered placement and \(d\)-layer base trajectories of \(A\) in the experiment above. Conditional on \(E\), denote by \(\mu_E\) the law of the overlay multiplier on \(A\) generated by the internal extra switches. We visualize its generation with labels starting at the corresponding indices of \(A\) (identity overlay at the start for describing this multiplier); in locations, these labels follow the actual switched paths. An internal encounter of such a label is when it is on a location pair having two indices in \(A\). The set of locations occupied by the labels from \(A\) is precisely the base location set of \(A\) throughout. Thus in the actual switch dynamics an internal encounter is pairing with another label from \(A\).

Fix constants
\[
 h=1/32,\qquad \eta=1/4096,\qquad \delta=\eta/4 .
\]
Call the extra-switch outcome good if at most \(q^{1-\eta}/2\) of its labels have fewer than \(h d\) internal encounters. Denote by \(\mu'_E\) the subprobability obtained from \(\mu_E\) by counting only switch outcomes that are good (there is no renormalization).

The mass discarded has expectation \(o(1)\). For every fixed starting placement, averaging over the base and extra switches gives the same product law of independent fair actual switch arrays, by predictability and the base-coin xor argument. Hence the actual location-switch array is independent of the uniformly sampled starting placement of \(A\), even though the choice of extra edges depends on that placement. In \(d\) consecutive directions, any two actual labels (considering labels on all the start sites) can be paired at most once. They could pair only at the direction of the \emph{last}, in layer order, of the bits on which their start sites differed. Before that direction such an unchanged different bit would preclude pairing on another direction. Afterwards all still unprocessed bits are equal, so pairing in an unprocessed direction (requiring a difference there) is impossible. Each label therefore has \(d\) distinct partners. Conditional on a specified active label's start position and on the actual location switch arrays, its partners' start positions are determined; the other \(q-1\) active start positions are a uniform subset among the \(n-1\) others. Its number of encounters is hypergeometric from sampling \(d\) sites out of \(n-1\) with \(q-1\) active. This law differs from a binomial with parameters \(d,(q-1)/(n-1)\) in total variation by at most \(\binom{d}{2}/(n-1)\), by coupling sampling with and without replacement. For large \(n\) the binomial mean is at least \(d/16\); the probability to be below \(hd\) is at most \(\exp(-d/128)\) by the binomial Chernoff bound. With the sampling error included, the resulting bound times \(q^\eta\) tends to zero. The expected number of low-encounter labels and Markov's inequality now give the claim on discarded mass.

Let \(\alpha,\beta\) be outcomes of independent extra-switch arrays given \(E\), identified also with their permutations. Write \(F\) for the number of fixed points of \(\alpha^{-1}\beta\). For all sufficiently large \(n\) and integers \(f\ge q^{1-\eta}\), we have
\begin{equation}
 \Pr(\alpha\text{ good},\ F\ge f)\ \le\ q^{-c_* f},
 \qquad c_*=h/8,\qquad (f\le q).                                      \label{special:overlay:eq:1}
\end{equation}
Indeed fix \(E\) and the switches of a good \(\alpha\). Agreement of the two paths of a label at their endpoint (i.e. agreement of the output index) requires agreement at every layer: in locations in the cyclic-direction frame, the start and end of a \(d\)-layer path specify it since each bit is processed only once. For any set of \(f\) agreed labels at least \(f/2\) have at least \(hd\) encounters in \(\alpha\)'s outcome. Making \(\beta\) follow those paths fixes its extra coins on their internal encounters, at least \(hdf/4\) distinct coins since each pair can be counted at most twice in its layer. These coins are independent given \(E\). Union bounding over sets costs at most \(\binom{q}{f}\le(e q/f)^f\); this times \(2^{-hdf/4}\) is at most \(q^{-h f/8}\) by \(d\ge\log_2 q\) and the choice of \(\eta\), for large \(n\). This proves \eqref{special:overlay:eq:1}.

We use irreducible complex unitary representations \(\pi_\lambda\) of \(\mathfrak S(A)\) with dimensions \(D_\lambda\) and (un-normalized) characters \(\chi_\lambda\). Write \(\widehat\nu(\lambda)=\sum_g\nu(g)\pi_\lambda(g)\) for the matrix of a (possibly subprobability) law. Define the per-segment factor
\begin{equation}
 b_\lambda=\frac{1}{D_\lambda}\mathbb E_E\operatorname{Tr}\!\left(
       \widehat{\mu'_E}(\lambda)^*\,\widehat{\mu'_E}(\lambda)\right)
  =\mathbb E\!\left[
       \mathbf 1_{\alpha\ {\rm good},\,\beta\ {\rm good}}\,
       \frac{\chi_\lambda(\alpha^{-1}\beta)}{D_\lambda}\right].          \label{special:overlay:eq:2}
\end{equation}
The second equality uses the independent switches of \(\alpha,\beta\) conditional on \(E\). Moreover the expectation of the matrix inside the trace is \(b_\lambda I\). The law of the subprobability kernel is invariant under conjugating by permutations of \(A\): this renames the uniformly placed indices, with the same good-outcome rule. The scalar matrix claim follows by Schur's lemma.
The representation-theoretic background used here is recalled in~\cite{sagan,etingof}.

Apply Lemma~\ref{lem:character-fixed-points} with group size \(q\) and \(\delta=\eta/4=1/16384\). Since \(4\delta=\eta\), if the number \(f\) of fixed points is at most \(q^{1-\eta}\), then
\[
 \frac{|\chi_\lambda(\sigma)|}{D_\lambda}\le D_\lambda^{-\delta}.
\]
For all \(f\), the same lemma gives
\begin{equation}
 \frac{|\chi_\lambda(\sigma)|}{D_\lambda}
      \le D_\lambda^{-\delta}q^{\delta f}.\label{special:overlay:eq:3}
\end{equation}

Applying the small-\(f\) bound and \eqref{special:overlay:eq:1}, \eqref{special:overlay:eq:3} in \eqref{special:overlay:eq:2} (bounding by absolute values), and using \(\delta<c_*\), yields
\begin{equation}
 0\le b_\lambda\ \le\ 2D_\lambda^{-\delta}                             \label{special:overlay:eq:4}
\end{equation}
for large \(q\), since \(\sum_{f\ge q^{1-\eta}} q^{-(c_*-\delta) f}=o(1)\). For the sign representation we have the better bound \(b_{\rm sign}=o(1)\). If the environment contains any internal match then the sign expectation before truncation is zero by the fair extra coin on it. Its magnitude after truncation is then at most the mass discarded. If there is no internal match, no outcome is good, for large \(q\). This proves the sign estimate by the bound on expected discarded mass.

The reciprocal-degree estimate of Lemma~\ref{lem:degree-sum} gives
\begin{equation}
 \sum_{\lambda:\,D_\lambda>1}D_\lambda^{-2}=o(1).
 \label{special:overlay:eq:5}
\end{equation}

Now take \(k=\lceil 4/\delta\rceil\). In the independent segment experiment let \(\nu'\) denote the product law obtained by convolving the \(k\) subprobabilities (multipliers applied on the left). At each nontrivial representation the expected squared Hilbert-Schmidt norm of its matrix is \(D_\lambda\) times the product of the \(b_\lambda\) factors: iteratively use independence and the scalar matrix expectation after \eqref{special:overlay:eq:2} for the latest matrix in the product. The estimates are uniform even if cyclic phases differ. By finite group Plancherel, for the uniform probability \(U_A\) on \(\mathfrak S(A)\),
\[
 q!\sum_g |\nu'(g)-U_A(g)|^2
    = |\nu'(\mathfrak S(A))-1|^2+
      \sum_{\lambda\ne{\rm triv}}D_\lambda
             \|\widehat{\nu'}(\lambda)\|_{\rm HS}^2 .
\]
The mass deficit tends to zero in expectation (bounded by the sum of discarded masses), thus also in expected square. The expected sign term tends to zero. The expected sum for \(D_\lambda>1\) is at most
\[
 2^k \sum_{D_\lambda>1} D_\lambda^{2-\delta k}=o(1)
\]
by \eqref{special:overlay:eq:4}, \eqref{special:overlay:eq:5}. These account for all irreducibles (the one-dimensional ones being trivial and sign for large \(q\)). By Cauchy-Schwarz the expected \(\ell^1\) error tends to zero. Reinstating the discarded masses costs \(o(1)\) also in expected \(\ell^1\), since the product law before truncation dominates \(\nu'\) pointwise. This proves Proposition~\ref{prop:overlay-call} for independent segments, and hence for the actual segments by the comparison in the call construction.

\subsection{From group calls to the whole permutation}

We can now randomize any one prescribed group of $n/8$ indices, conditional on the base in expected total variation. Mixing the groups of just one partition would preserve which labels belong to each group. Two transverse partitions remove that obstruction. We prove the required three-round comparison, including its discrete total-variation conclusion.

Choose two partitions of the base index set \(V\), one into \(R_1,\ldots,R_K\), the other into \(C_1,\ldots,C_K\), with \(|R_i\cap C_j|=n/K^2\); this is possible for large \(d\). Make calls for the \(R\)-groups, then the \(C\)-groups, then the \(R\)-groups again, one call for each group of the indicated partition, in any fixed order within each partition. Conditional on the full base, the multipliers from distinct calls are independent (they use independent extra coins on edges prescribed from the base). In expectation over the base, their product law has total-variation error \(o(1)\) from the product law using independent uniforms on the called groups. This follows by the per-call error just proved, summed over the \(3K\) calls, replacing conditional laws in the product and mapping to the composed permutation. Inside each partition these ideal uniforms together act as an independent uniform shuffling within each group.

The remaining comparison concerns only uniform group shuffles. Alternating independent row and column permutations were studied by H\aa stad in the square lattice shuffle~\cite{hastad2006}. Here the number of blocks is fixed and their intersections grow; the proof uses the resulting contingency-table local limit.

\begin{lemma}[Three rounds on transverse partitions]\label{lem:transverse-rounds}
Let \(R_1,\ldots,R_K\) and \(C_1,\ldots,C_K\) partition \(V\), with \(|R_i\cap C_j|=n/K^2\). Independently shuffle uniformly within every \(R\)-group, then within every \(C\)-group, and then within every \(R\)-group again. For the fixed \(K=8\) and \(n=2^d\), the resulting permutation law has total-variation distance \(o(1)\) from uniform on \(\mathfrak S(V)\) as \(d\to\infty\).
\end{lemma}
\begin{proof}
The product law is invariant on both right and left by the product subgroup of permutations within the \(R\)-groups, due to the first and last ideal shuffles. As with a uniform permutation, its probabilities are constant on permutations having the same count matrix \(X=(X_{ij})\), where \(X_{ij}\) counts moves from \(R_i\) to \(R_j\). Indeed permutations with equal counts can be identified by rearranging inputs within the groups (matching their output groups) and then rearranging outputs within the groups. Thus it suffices to compare just the count laws.

For a uniform permutation let this table law be \(p_n\); in general write \(p_M\) for the analogous law on \(M\) items for \(K\) groups of equal size \(m=M/K\). Its probability is zero off nonnegative integer tables with all row and column sums \(m\), and on such tables is
\begin{equation}
 p_M(x)=\frac{(m!)^{2K}}{M!\prod_{i,j}x_{ij}!}.                        \label{special:overlay:eq:6}
\end{equation}
This follows by partitioning each input group into its outputs' groups, then ordering the inputs assigned to each output group. In the ideal shuffle, at the middle stage each \(C\)-group contains \(n/K^2\) labels from each initial \(R\)-group. Within each \(C\)-group these labels are uniformly permuted with also \(n/K^2\) slots belonging to each \(R\)-group. Conditional on the first shuffles the resulting tables within the \(C\)-groups are independent with law \(p_{n/K}\), independent in law of that conditioning. The last \(R\)-shuffles do not change total counts. Write \(T^{(\ell)}\) for the row-to-row count table contributed by \(C_\ell\). These tables are independent with law \(p_{n/K}\), each has row and column sums \(a_0=n/K^2\), and \(X_{ij}=\sum_{\ell=1}^K T^{(\ell)}_{ij}\). Figure~\ref{fig:overlay-grid} distinguishes the geometric columns from the destination-group columns of these tables.

\begin{figure}[htbp]
\centering
\begin{tikzpicture}[
  x=1cm,y=1cm,>=Latex,
  every node/.style={font=\small},
  gridline/.style={draw=black!45,line width=.35pt},
  tableentry/.style={font=\scriptsize},
  note/.style={font=\footnotesize}
]
  \foreach \row/\tone in {1/blue!12,2/teal!16,3/orange!18,4/violet!13} {
    \pgfmathsetmacro{\yb}{(4-\row)*.68}
    \fill[\tone] (0,\yb) rectangle (2.72,\yb+.68);
    \node[anchor=east] at (-.14,\yb+.34) {$R_{\row}$};
    \foreach \col in {1,...,4} {
      \pgfmathsetmacro{\xc}{(\col-.5)*.68}
      \node at (\xc,\yb+.34) {$a_0$};
    }
  }
  \foreach \line in {0,...,4} {
    \draw[gridline] (\line*.68,0)--(\line*.68,2.72);
    \draw[gridline] (0,\line*.68)--(2.72,\line*.68);
  }
  \foreach \col in {1,...,4} {
    \node at ({(\col-.5)*.68},2.98) {$C_{\col}$};
  }
  \draw[black!85,line width=1.1pt] (2.04,0) rectangle (2.72,2.72);
  \node at (1.36,3.48) {Before the column shuffles};
  \node[note,align=center] at (1.36,-.51)
    {Each cell contains $a_0=n/K^2$ labels\\of its initial row color.};

  \draw[->,line width=.7pt] (2.79,1.36)--(4.80,1.36);
  \node[note,align=center] at (3.79,1.90) {shuffle within\\one $C_\ell$};
  \node[note] at (3.79,.97) {(here $\ell=4$)};

  \foreach \row in {1,...,4} {
    \pgfmathsetmacro{\yb}{(4-\row)*.68}
    \node[anchor=east] at (5.84,\yb+.34) {$R_{\row}$};
    \foreach \col in {1,...,4} {
      \pgfmathsetmacro{\xc}{6.00+(\col-.5)*.68}
      \node[tableentry] at (\xc,\yb+.34) {$T^{(\ell)}_{\row\col}$};
    }
    \node[note] at (9.00,\yb+.34) {$a_0$};
  }
  \foreach \line in {0,...,4} {
    \draw[gridline] (6.00+\line*.68,0)--(6.00+\line*.68,2.72);
    \draw[gridline] (6.00,\line*.68)--(8.72,\line*.68);
  }
  \foreach \col in {1,...,4} {
    \node at ({6.00+(\col-.5)*.68},2.98) {$R_{\col}$};
    \node[note] at ({6.00+(\col-.5)*.68},-.23) {$a_0$};
  }
  \node[rotate=90,note] at (5.13,1.36) {initial $R_i$};
  \node[note] at (7.36,3.48) {final $R_j$};
  \node[note] at (7.36,-.62) {Row and column sums are $a_0$.};

  \node[align=center] at (4.40,-1.40)
    {$T^{(1)},\ldots,T^{(K)}$ independent,\quad
      $T^{(\ell)}\sim p_{n/K}$,\qquad
      $X_{ij}=\displaystyle\sum_{\ell=1}^{K}T^{(\ell)}_{ij}$.};
\end{tikzpicture}
\caption{The middle round in the row--column--row construction
(schematic $K=4$; the proof uses $K=8$).
Every $C_\ell$ contains $a_0$ labels from each initial $R_i$ and $a_0$ slots
in each final $R_j$. Independent column shuffles produce the tables
$T^{(\ell)}$; their columns index final $R_j$ groups.
The first and last row shuffles supply right and left invariance,
respectively, making the permutation law uniform conditional on its count matrix.}
\label{fig:overlay-grid}
\end{figure}

Here is the comparison of this sum with \(p_n\). We may assume \(n\) is divisible by \(K^3\). Use as free coordinates the upper-left \((K-1)\times(K-1)\) subtable, treated as a vector of size \(b=(K-1)^2\). For \(M\) divisible by \(K^2\) and \(X\sim p_M\), center the vector at its coordinate mean \(M/K^2\) and divide by \(\sqrt M\); call the result \(Z_M\). Let \(J\) be the linear map completing a free-coordinate vector to a full matrix with zero row and column sums. Uniformly for \(\|z\|\) bounded on the centered scaled lattice,
\begin{equation}
 \Pr(Z_M=z)=M^{-b/2}
   \left(G(z)+o(1)\right),\qquad
 G(z)=c_K \exp\!\left(-K^2\|J z\|^2/2\right),                           \label{special:overlay:eq:7}
\end{equation}
where matrix norm here is Euclidean and \(c_K>0\) depends only on \(K\). To see this, substitute \(x_{ij}=a+\sqrt M (Jz)_{ij}\) with \(a=M/K^2\) in \eqref{special:overlay:eq:6} (nonnegativity holds for bounded \(z\) and large \(M\)). Stirling's formula and expansion of \(x\log x-x\) about \(a\) have canceling linear terms in the sum. The constant terms cancel with the numerator and \(M!\) and the quadratic terms give \(K^2\|Jz\|^2/2\), with remainders \(o(1)\) for bounded \(z\). The square-root factors give \(M^{K-1/2-K^2/2}=M^{-b/2}\) times a constant \((2\pi)^{-b/2}K^{K^2-K}\) and \(1+o(1)\), proving \eqref{special:overlay:eq:7}.

Writing \(W\) for an independent uniform vector on \([0,1)^b\), \(Z_M+W/\sqrt M\) has density converging in \(L^1\) to \(G\). On compacts this follows from \eqref{special:overlay:eq:7} and continuity. Tail probabilities are uniformly small on taking large compacts since each count coordinate is hypergeometric with variance at most \(M\); the limit density is integrable since \(J\) is injective. This proves the \(L^1\) convergence and in particular normalization, so \(G\) is a centered nondegenerate Gaussian density.

For the sum of \(K\) tables from \(p_{n/K}\), the centered vector at scale \(1/\sqrt n\) is a sum of \(K\) independent \(Z_{n/K}/\sqrt K\). First add independent noise \(W_i/\sqrt n\) to each summand. By the density convergence, independence and total-variation contraction under addition, the noisy sum converges in total variation to the law with density \(G\): the sum of \(K\) centered Gaussians of this law each divided by \(\sqrt K\) has the same law. This convergence still holds keeping only the noise on the first summand. Indeed that smoothed summand by itself converges in total variation to a Gaussian (scaled by \(1/\sqrt K\)). Its distance to its translate by any vector of norm at most \((K-1)\sqrt b/\sqrt n\) therefore tends to zero uniformly, by continuity of translations of the Gaussian density in \(L^1\). Conditioning on the other summands and their noises bounds by this \(o(1)\) the distance induced by removing their noises. Thus, with a single noise, both the centered scaled sum and the \(p_n\) vector are close to the same law by \eqref{special:overlay:eq:7} and the sum argument. Total variation between these two noisy laws equals the discrete total variation without noise since both lattices have the same cells spread out uniformly by the noise. This proves the desired count-law comparison, hence mixing of the ideal product.
\end{proof}

\label{special:overlay:bound}
We conclude that the overlay product conditional on the base has expected total-variation distance \(o(1)\) from a uniform permutation on \(V\). Multiplying by the final base permutation (fixed under the conditioning) preserves the comparison with uniform. Removing the conditioning and mapping back to the deck locations proves total-variation distance \(o(1)\) for the actual shuffle. All estimates hold for any initial deck by identifying its cards with their starting positions. The shuffle has uniform stationary law since every layer applies a random permutation according to switches independent of the incoming deck. We used \(3K\cdot k\cdot 21d\) full layers of the type counted as steps in the shuffle definition. This gives the upper bound by an absolute multiple of \(d\) for sufficiently large \(d\).
The choices give $\delta=1/16384$, $k=65536$, and $3Kk\cdot21=33030144$, so the upper time is $33030144d$ complete shuffles.

This proves Theorem~\ref{thm:overlay}. The support bound of Section~\ref{sec:conditional} also gives $t_{\mathrm{mix}}\ge d$ for all sufficiently large $d$.
